\documentclass[10pt,a4paper]{amsart}
\usepackage[margin=19mm]{geometry}
\usepackage{amsmath,amssymb,mathtools}
\usepackage[scr=rsfs,cal=euler]{mathalfa}
\usepackage{xfrac}
\usepackage[unicode,hidelinks,bookmarks=false]{hyperref}
\newcommand{\ie}{i.e.\ }
\newcommand{\cf}{cf.\ }
\newcommand{\resp}{respectively\ }
\newcommand{\Subsection}{\subsection}
\providecommand{\nobreakdash}{\nobreak\hbox{-}}
\setlength{\emergencystretch}{2em}
\multlinegap0pt
\usepackage{bm}
\usepackage{bbm}
\usepackage{dsfont}
\usepackage{xspace}
\usepackage{cancel}
\usepackage{mathtools}
\usepackage{amsthm}
\makeatletter
\@ifundefined{definition}{\newtheorem{definition}{Definition}}{}
\@ifundefined{lemma}{\newtheorem{lemma}{Lemma}}{}
\@ifundefined{theorem}{\newtheorem{theorem}{Theorem}}{}
\makeatother
\makeatletter
\expandafter\ifx\csname note\endcsname\relax
\newenvironment{note}{{\it Note:}}{}
\fi
\expandafter\ifx\csname proof\endcsname\relax
\newenvironment{proof}{{\noindent\it Proof:} }{\hfill $\square$\par}
\fi
\makeatother
\title[Symplectic self-orthogonal quasi-cyclic codes]{Symplectic self-orthogonal quasi-cyclic codes}
\author{Chaofeng Guan, Ruihu Li, Jingjie Lv, Zhi Ma}
\date{Source: arXiv:2212.14225v5 (CC BY 4.0)}
\begin{document}
\maketitle

\noindent\emph{Source and licence.} Symplectic self-orthogonal quasi-cyclic codes. Version arXiv:2212.14225v5 (CC BY 4.0), distributed under the Creative Commons Attribution 4.0 International licence, as recorded in the arXiv licence metadata for this exact version and shown on the abstract page https://arxiv.org/abs/2212.14225v5. Full rights notice: licenses/arxiv-2212.14225-CC-BY-4.0.txt.\par\smallskip

\noindent\textit{Editorial scope.} Counted unit: the Theorem whose source label is \texttt{Dual\_Structure} in the article (source0.tex lines 784--798, source label \texttt{Dual\_Structure}), delivered together with the article's own complete original proof of it (source0.tex lines 799--821, one \texttt{proof} environment of 23 lines). No mathematics is written, repaired or re-derived: statement and proof are the article's own text line for line. Required context retained uncounted: one\_quasi-cyclic\_def (lines 300--309); one\_quasi-cyclic (lines 314--321); QCmatrix\_spread (lines 532--538). Every definition, display and label of those lines is retained unchanged. Omitted from this package and not redistributed: 1--299, 310--313, 322--531, 539--783, 822--1338. Nothing inside a retained excerpt is omitted or altered: every retained body is delivered exactly as the article's lines are written, so no omission receipt of any kind accompanies this package. Citations: the retained excerpts carry no citation command at all, so no bibliography entry of the article is transcribed and no reference is redistributed. Reference adaptation: none was needed and none was made; every reference inside the delivered unit resolves inside it, so no unresolved reference or citation is produced. Printed slips kept literal: no printed word, symbol, hypothesis or number of the counted statement or of its proof was altered. Layout and class: the article's own class is \texttt{IEEEtran} with packages including amsmath, amsfonts, algorithmic, algorithm, array, subfig, textcomp, stfloats, url, verbatim, graphicx, amssymb, cuted, mathrsfs; the wrapper is the lane's pinned \texttt{amsart} editorial wrapper at 10pt on A4, which restates the theorem-like environments the retained text needs and adds the article's own macro definitions that the retained text needs (none), with extra wrapper packages (bm, bbm, dsfont, xspace, cancel, mathtools). The delivered numbering is the unit's own. Assets: no retained span contains a figure environment, a graphics-inclusion command, a PDF-page inclusion, a table or a TikZ picture; no image file is copied into this package, the package declares no asset dependency and no image file is redistributed with it. Licence: version arXiv:2212.14225v5 of this article is distributed under the Creative Commons Attribution 4.0 International licence, as recorded in the arXiv licence metadata for this exact version and shown on the abstract page https://arxiv.org/abs/2212.14225v5. A full rights notice is delivered at licenses/arxiv-2212.14225-CC-BY-4.0.txt. Characters: every retained excerpt is pure ASCII, so the delivered engine, XeLaTeX with its default Latin Modern text font, reports no missing character.
\begin{definition}\label{one_quasi-cyclic_def}
	Let $g(x)$ be a polynomial in $\mathbb{R}_{q,n}$ that divides ${{x}^{n}}-1$, and let $f_{j}(x)$ be $\ell$ different polynomials in $\mathbb{R}_{q,n}$ for $0\le j\le \ell-1$, such that the greatest common divisor of $f_0(x),f_1(x),\ldots,f_{\ell-1}(x)$ and $\frac{x^n-1}{g(x)}$ is equal to $1$. 
	%Let $g(x)$ and ${{f}_{j}}(x)$ be polynomials in $\mathbb{R}_{q,n}$, where $g(x)\mid ({{x}^{n}}-1)$, and let $h(x)=\frac{x^n-1}{g(x)}$, $\gcd(f_0(x),f_1(x),\ldots,f_{\ell -1}(x),h(x))=1$. 
	If $\mathscr{C}$ is a QC code generated by $(\left[g(x){f}_{0}(x)\right]$, $\left[g(x){f}_{1}(x)\right],\ldots,\left[g(x){f}_{\ell -1}(x)\right]) \in \mathbb{R}_{q,n}^{\ell} $, then $\mathscr{C}$ is called a $1$-generator QC code with index $\ell$.
	The generator matrix $G$ of $\mathscr{C}$  has the following form:
	\begin{equation}
		G=\left(G_{0}, G_{1}, \ldots, G_{\ell-1}\right) ,
	\end{equation}
	where $G_ {j} $ is an $n\times n$  circulant matrix generated by $\left[g (x) f_{j} (x)\right]$.
\end{definition}

\begin{theorem}\label{one_quasi-cyclic}
	If $\mathscr{C}$ is a $1$-generator QC code in Definition \ref{one_quasi-cyclic_def} with index even $\ell$, then $\mathscr{C}$ is symplectic self-orthogonal if and only if the following equation holds:
	\begin{equation}
		%{{g}^{{{\bot}_{e}}}}(x)\mid \sum\limits_{j=0}^{m-1}{g}(x)({{f}_{j}}(x){{\bar{f}}_{m+j}}(x)-{{f}_{m+j}}(x){{\bar{f}}_{j}}(x)).
		g^{\perp_e}(x)\mid g(x)\Lambda_1,
	\end{equation}
where $\Lambda_1 =\sum\limits_{j = 0}^{m - 1}( {f_j}(x){{\bar f}_{m + j}}(x) - {f_{m + j}}(x){{\bar f}_j}(x) )$ and $m=\frac{\ell}{2}$.
\end{theorem}

\begin{lemma}\label{QCmatrix_spread}
	With the above notation.
	Let $\mathscr{C}$ be a $1$-generator QC code with generator $(\left[f_0(x)g(x)\right],\left[f_1(x)g(x)\right])$, and let $t=n-\deg(g(x))-\deg(\gcd( h(x),f_1(x)))-\deg(\gcd( h(x),f_0(x)))$.
	Suppose that $\mathscr{C}_1$, $\mathscr{C}_2$ and $\mathscr{C}_3$ are codes generated by $\left[\frac{x^n-1}{\gcd( h(x),f_1(x))}\right]$, $\left[ \frac{x^n-1}{\gcd( h(x),f_0(x))}\right]$ and $\mathrm{Circ}_t\left( \left[g(x)f_0(x)\right],\left[g(x)f_1(x)\right] \right)$, where $h(x)=\frac{x^n-1}{g(x)}$. Then,  $\mathscr{C}=\mathscr{C}_1\oplus \mathscr{C}_2 + \mathscr{C}_3$.
%	If $g(x)f_0(x)$ and $g(x)f_1(x)$ satisfy $g(x)f_0(x)\mid \frac{x^n-1}{\gcd(h(x),f_1(x))} \pmod{x^n-1}$ and $g(x)f_0(x)\mid \frac{x^n-1}{\gcd(h(x),f_1(x))} \pmod{x^n-1}$ .
	%there exist polynomials $r_1(x)$,$r_2(x)$ and $r_3(x)$ in $\mathbb{R}_{q,n}$ such that a generator matrix $G$ of the 
\end{lemma}

\begin{theorem}\label{Dual_Structure} 
	With the above notation,
	let $\mathscr{C}$ be a $1$-generator QC code with generator $(\left[f_0(x)g(x)\right],\left[f_1(x)g(x)\right])$.
	 If $\gcd(f_0(x),g(x))=1$,		
then the symplectic dual code $\mathscr{C}^{\perp_s}$ of $\mathscr{C}$ is a 2-generator QC code, whose generator matrix is
\begin{equation}
	G^{\perp_{s}}=
	\begin{pmatrix}
		F_{0}	& F_{1}\\
		\mathbf{0}&G_{g^{\perp_e}}
	\end{pmatrix},
\end{equation}
where $(F_{0},F_1)$ and $(\mathbf{0},G_{g^{\perp_e}})$ are the generator matrices of $1$-generator QC codes determined by $(\left[ \bar{f}_{0}(x)\right], \left[ \bar{f}_{1}(x)\right])$ and $(\left[ 0\right], \left[ g^{\perp_e}(x)\right])$, respectively. 
 %where $(F_{0},F_1)$ is an $n \times 2n$ circulant matrices generated by $(\left[ \bar{f}_{0}(x)\right], \left[\bar{f}_{1}(x)\right])$, $G_{g^{\perp_e}}$ is the generator matrix of the cyclic code generated by $\left[g^{\perp_e}(x)\right]$, and $\mathbf{0}$ is an appropriate zero matrix.
\end{theorem}

\begin{proof}
	 Let $\mathscr{C}^{\circ}$ denote the code generated by $G^{\perp_{s}}$.  
 Then any codeword of $\mathscr{C}$ and $\mathscr{C}^{\circ}$ has the form $\mathbf{c_{1}}=(\left[a(x) f_{0}(x) g(x)\right],$ $\left[a(x) f_{1}(x) g(x)\right])$ and $\mathbf{c_{2}}=(\left[b(x) \bar{f}_{0}(x)\right],\left[b(x) \bar{f}_{1}(x)+c(x)g^{\perp_e}(x)\right])$, where $a(x),b(x),c(x)\in\mathbb{R}_{q,n}$. 
	Using a similar approach of Theorem \ref{one_quasi-cyclic}, it is easy to prove that $\left\langle \mathbf{c_{1}}, \mathbf{c_{2}}\right\rangle_{s}=\mathbf{c_{1}} \cdot \Omega _{mn} \cdot \mathbf{c_{2}}^{T}=0$.
	Therefore, $\mathscr{C}$ and $\mathscr{C}^{\circ}$ are mutually symplectic orthogonal.
%	According to Definition \ref{one_quasi-cyclic_def}, $\mathscr{C}$ is a $1$-generator quasi-cyclic code has generator $(\left[g(x){f}_{0}(x)\right]$, $\left[g(x){f}_{1}(x)\right]$,$\ldots$, $\left[g(x){f}_{\ell -1}(x)\right])$, where $0\le j\le \ell -1$. 

 		%Since $\gcd(\bar{f}_{m-1}(x),x^n-1)=1$, the rank of $(F_{m-1}, F_{j})$ is $n$. In addition, 

 		By Lemma \ref{QCmatrix_spread}, the generator matrix of $\mathscr{C}^{\circ}$ also can be transformed into 
\begin{equation}
	G_1^{\perp_{s}}=
	\begin{pmatrix}
		F_{0}^{*}	& F_{1}^{*}\\
		A	& \mathbf{0}\\
		\mathbf{0}	& B\\
		\mathbf{0}&G_{g^{\perp_e}}
	\end{pmatrix},
\end{equation}
 		where $(F_{0}^{*},F_1^{*})$ denotes the first $(n-\deg(\gcd(x^n-1,\bar{f}_1(x)))-\deg(\gcd(x^n-1,\bar{f}_{0}(x))))$-row of $(F_{0},F_1)$, $A$ and $B$ are the generator matrices of cyclic codes generated by $\left[  \frac{x^n-1}{\gcd(x^n-1,\bar{f}_{1}(x))}\right]$ and $\left[ \frac{x^n-1}{\gcd(x^n-1,\bar{f}_{0}(x))}\right]$, respectively.
 		
 		According to Definition \ref{one_quasi-cyclic_def}, $\gcd(\bar{f}_{0}(x),\bar{f}_{1}(x),g^{\perp_e}(x))=1$. Given that $\gcd(f_0(x),g(x))=1$, we have $\gcd(\bar{f}_0(x),\tilde{g}(x) )=1$. As $\tilde{g}(x)g^{\perp_e}(x)=x^n-1$,  $\gcd(\bar{f}_{0}(x),\bar{f}_{1}(x),x^n-1)=1$ and $\operatorname{lcm}\left(g^{\perp_e}(x),\frac{x^n-1}{\gcd(x^n-1,\bar{f}_{0}(x))}\right) \equiv 0 \pmod{x^n-1}$ hold. This implies that, the rank of $(F_{0}, F_{1})$ is $n$ and the rows of $B$ and $G_{g^{\perp_e}}$ are linearly independent. Therefore, the rank of $G_1^{\perp_{s}}$ is $n+\deg(g(x))$. Thus, $\mathscr{C}$ and $\mathscr{C}^{\circ}$ are symplectic dual codes.
\end{proof}

\end{document}
